\documentclass[11pt,a4paper]{article}
\usepackage{../shared/luxcover}
\usepackage[utf8]{inputenc}
\usepackage[T1]{fontenc}
\usepackage{amsmath,amssymb,amsfonts,amsthm}
\usepackage{graphicx}
\usepackage{hyperref}
\usepackage{booktabs}
\usepackage{enumitem}
\usepackage{xcolor}
\usepackage{siunitx}
\usepackage{listings}
\usepackage{longtable}
\usepackage{array}
\input{../shared/lstlang}

\hypersetup{colorlinks=true,linkcolor=black,citecolor=blue,urlcolor=blue}

\lstset{
  basicstyle=\ttfamily\scriptsize,
  breaklines=true,
  frame=single,
  showstringspaces=false,
  numbers=left,
  numberstyle=\tiny\color{gray},
  keywordstyle=\color{blue!70!black},
  commentstyle=\color{green!45!black},
  stringstyle=\color{red!60!black},
}

% Provenance tags — every quantitative claim carries one (house convention, matches
% gpu-native-billion-tps.tex).
\newcommand{\Mnow}{\textbf{\textcolor{black}{M}}}   % measured by us this session (2026-07-10)
\newcommand{\Mprior}{M$^\star$}                      % measured in prior cited work
\newcommand{\Proj}{P}                                % projected / derived by calculation
\newcommand{\Claim}{C}                               % external / cited claim

\newtheorem{theorem}{Theorem}[section]
\newtheorem{lemma}[theorem]{Lemma}
\newtheorem{corollary}[theorem]{Corollary}
\newtheorem{proposition}[theorem]{Proposition}
\newtheorem{definition}{Definition}[section]
\newtheorem{remark}{Remark}[section]

\title{Nova: A Two-Tier Metastable Finality Ladder for Lux\\[4pt]
       \large Majority-accept and \(\tfrac{2}{3}\)-stake export, decomposed --- the v1.36 consensus}
\author{Lux Industries, Inc. \\
        \texttt{research@lux.network}}
\date{2026-07-10}

\begin{document}
\luxcoverpage

\begin{abstract}
Lux v1.36 (\texttt{consensus v1.36.0}, \texttt{node v1.36.0}, \texttt{evm v1.104.8})
ships \emph{Nova}: a leaderless, post-quantum, metastable consensus whose defining
structural choice is a \textbf{two-tier finality ladder}. A block climbs five distinct
authorization rungs --- \(\textbf{Photon} \to \textbf{Wave} \to \textbf{Nova} \to
\textbf{Quasar} \to \textbf{Horizon}\) --- and the two load-bearing rungs are kept
\emph{orthogonal}. \textbf{Nova} is \emph{local acceptance} at a simple majority of the
validator set, \(\mathrm{NovaQuorum}(n)=\lfloor n/2\rfloor + 1\); it is crash-fault-safe
and authorizes local execution but is reorgable. \textbf{Quasar} is \emph{export
finality}: a strict \(>\tfrac{2}{3}\)-of-stake, Byzantine-safe post-quantum certificate
that alone authorizes bridges, DEX settlement, and cross-chain messages. The invariant is
one line: \emph{a block below Quasar can never be exported}. This separation --- majority
to make local progress, supermajority to make irreversible external claims --- is why a
five-validator fleet keeps producing at three live nodes (where a single
\(\tfrac{2}{3}\)-only engine freezes) while never exporting an unsafe block.

We specify the ladder against the shipped code (\texttt{engine/chain/finality.go},
\texttt{quorum.go}, \texttt{cert.go}, \texttt{config/quorum\_threshold.go}), give the
safety arguments (crash-fault majority-intersection for Nova; the Byzantine overlap bound
\(2\alpha - K > f\) and a no-two-conflicting-certificates theorem for Quasar), describe the
ZAP-native plugin transport (\texttt{luxfi/zap}, no gRPC, no BFT round machine), and show
the rule is size-invariant from \(n=1\) to \(n=10^{6}\) because \(K\) is a \emph{sample}
size, not the set size. We then report measured evidence: a live \num{5}-validator testnet
(\texttt{ghcr.io/luxfi/node:v1.36.0}) held common-height coherence through a
\(5/5 \to 4/5 \to 3/5\) node-drop while producing (\(+20,+8,+6\) blocks per \SI{70}{s}
phase), and a steady load run measured the live C-Chain mining at
\(\sim\!\SI{1.6}{\percent}\) block fullness with the base fee pinned at its floor ---
locating the live throughput ceiling in the EVM fee/gas-target throttle, \emph{not} in
consensus, which stayed coherent and idle. Every number is tagged by provenance
(\Mnow\ measured-now, \Mprior\ measured-prior, \Proj\ projected, \Claim\ cited).
\end{abstract}

\tableofcontents
\newpage

\section{Introduction: one ladder, five rungs, two tiers that matter}

A consensus stack answers one hard question --- \emph{when is a block irreversible?} ---
and a well-built stack answers it in exactly one place. Lux answers it with a
\textbf{finality ladder}: an ordered set of authorization tiers, each a strictly higher
authority than the last, with the boundaries between them never collapsed. The ladder is
named in the Lux astrophysics ontology --- light is emitted, it propagates, a star
ignites, it blazes into the brightest beacon in the sky, it crosses the event horizon:

\begin{center}
\(\textbf{Photon} \;\to\; \textbf{Wave} \;\to\; \textbf{Nova} \;\to\; \textbf{Quasar}
\;\to\; \textbf{Horizon}\)
\end{center}

\begin{center}
\begin{tabular}{@{}l l l@{}}
\toprule
\textbf{Rung} & \textbf{Meaning} & \textbf{Authorizes} \\
\midrule
Photon  & a vote is in flight; the block is being sampled           & nothing \\
Wave    & the sampling wave has made this block the build tip       & nothing \\
\textbf{Nova}   & \(\beta\)-ignition: a \emph{majority} sustained \(\beta\) rounds & \textbf{local execution} (reorgable) \\
\textbf{Quasar} & a strict \(>\tfrac{2}{3}\)-stake certificate exists       & \textbf{export}: bridges / DEX / cross-chain \\
Horizon & a post-quantum signature seals the certificate            & irreversible settlement \\
\bottomrule
\end{tabular}
\end{center}

The two rungs that carry the protocol's weight are \textbf{Nova} and \textbf{Quasar}, and
the central design decision of v1.36 is that \emph{they are different tiers with different
thresholds and different jobs}:

\begin{itemize}[leftmargin=1.4em]
\item \textbf{Nova} is \emph{local acceptance}. It fires at a \emph{simple majority} of the
validator set. It is crash-fault-safe --- it survives any minority going down --- and it
authorizes the node to execute the block locally. It is \emph{reorgable}: a lone equivocator
can, in principle, split a bare majority, so Nova is safe to \emph{act on} but never to
\emph{export}.
\item \textbf{Quasar} is \emph{export finality}. It fires only when a strict
\(>\tfrac{2}{3}\) of \emph{stake} has signed --- the Byzantine-safe threshold --- and it is
the \emph{only} tier that may drive an irreversible external effect: a bridge withdrawal, a
\texttt{0x9999} DEX settlement receipt, a cross-chain message, a validator-set transition.
\end{itemize}

\begin{quote}
\textbf{THE INVARIANT.} A block below Quasar can \emph{never} be exported as
deterministically final to a bridge, another chain, or an irreversible settlement. Nova
authorizes local execution; export requires Quasar; quantum-irreversibility requires
Horizon. (\texttt{engine/chain/finality.go}.)
\end{quote}

This paper is the rigorous, measured treatment of that ladder. It is the companion to the
normative specification \textbf{LP-305}~\cite{lp305} and the cross-system landscape
\textbf{LP-306}~\cite{lp306}; where this paper and the code disagree, the code
(\texttt{consensus v1.36.0}) governs and this paper is the defect.

\paragraph{Why two tiers and not one.} A single threshold cannot be both maximally live and
maximally safe. Set the bar at a majority and you get liveness down to \(n/2\) but you
cannot make a Byzantine-safe external claim. Set it at \(\tfrac{2}{3}\) of stake and you get
a Byzantine-safe claim but you freeze the moment more than a third is unreachable --- a
five-node fleet stops at three live nodes. Nova refuses the false choice: it runs \emph{both}
thresholds, on \emph{different rungs}, for \emph{different purposes}. The majority tier keeps
the chain making local progress; the supermajority tier gates only what leaves the chain. A
node that has lost its supermajority keeps accepting and executing (Nova) while \emph{honestly}
declining to export (no Quasar) until the supermajority returns.

\paragraph{Provenance discipline.} This is a measurement paper; it does not launder
projections as results. Every quantitative claim is tagged: \Mnow{} = measured by us this
session (2026-07-10) on the hardware named in \S\ref{sec:eval}; \Mprior{} = measured in
prior cited Lux work; \Proj{} = projected or derived by calculation; \Claim{} = an external
cited claim. Untagged numbers are definitions.

\section{Nova --- majority local-acceptance}
\label{sec:nova}

Nova is the \emph{lower} of the two load-bearing tiers and the one that carries chain
liveness. It is driven by a leaderless metastable sampler (Photon committee sampling
\(\to\) Wave / FPC threshold voting \(\to\) Focus \(\beta\)-confidence; LP-111 / LP-027 /
LP-028 / LP-114), and it \emph{ignites} --- transitions a block to the Nova rung and calls
\texttt{VM.Accept} --- when a majority of the validator set has sustained the block for
\(\beta\) consecutive rounds.

\subsection{The quorum: \(\mathrm{NovaQuorum}(n)=\lfloor n/2\rfloor + 1\)}

The Nova ignition threshold is a clean function of the live validator count \(n\)
(\texttt{engine/chain/quorum.go}):

\begin{lstlisting}[language=Go]
// NovaQuorum is the majority the sampler needs to IGNITE a block to Nova
// (acceptance, local execution). It tolerates floor((n-1)/2) crash faults --
// the MOST a net of n can lose and still make safe local progress under the
// crash-fault model: any two majorities intersect, a non-equivocating node
// cannot back two conflicting blocks, and a partition lets at most one side
// hold a majority. It is DELIBERATELY below the Quasar 2/3 stake gate.
func NovaQuorum(n int) int {
    if n < 1 { return 1 }   // a lone node self-ignites; never 0
    return n/2 + 1
}

// CrashTolerance = n - NovaQuorum(n) = ceil(n/2) - 1: crash faults Nova
// ignition survives. Beyond it Nova PAUSES fail-closed (no ignition => no
// acceptance => no fork) and resumes the instant a node returns.
func CrashTolerance(n int) int {
    if n < 2 { return 0 }
    return n - NovaQuorum(n)
}
\end{lstlisting}

\begin{definition}[Nova quorum and crash tolerance]
For a validator set of size \(n\), \(\mathrm{NovaQuorum}(n)=\lfloor n/2\rfloor+1\) and
\(\mathrm{CrashTolerance}(n)=n-\mathrm{NovaQuorum}(n)=\lceil n/2\rceil-1\).
\end{definition}

For the live \(n=5\) primary-network fleet this is \(\mathrm{NovaQuorum}(5)=3\) and
\(\mathrm{CrashTolerance}(5)=2\): \textbf{Nova keeps accepting with any three of five nodes
responsive}. This single fact is the operational heart of v1.36 --- and it is exactly where a
\(\tfrac{2}{3}\)-only engine dies, because \(\lceil 2\cdot 5/3\rceil = 4 > 3\).

\subsection{The confidence depth \(\beta\)}

Ignition additionally requires \(\beta\) \emph{consecutive} majority rounds (hysteresis), so
a one-round flap cannot ignite (\texttt{quorum.go}):

\begin{lstlisting}[language=Go]
// NovaBeta: n==1 -> 1 (a lone node ignites immediately; any wait is a
// phantom-peer stall). 2<=n -> 2 (one confirming round of hysteresis, without
// the fragility of a deep beta at zero margin). On K=n nets beta is NOT a
// Byzantine knob -- that is Quasar's job; large permissionless nets (K<n, real
// subsampling) raise beta to the classical negl(beta) metastable depth.
func NovaBeta(n int) int { if n <= 1 { return 1 }; return 2 }
\end{lstlisting}

\subsection{Fail-closed, never fork-open}

When responsive stake drops below \(\mathrm{NovaQuorum}(n)\), Nova does not lower its bar and
it does not guess --- it \textbf{pauses}. No ignition means no acceptance means no fork; the
chain halts safely and resumes, via bounded event-driven re-poll and runtime ancestor
catch-up (LP-303~\cite{lp303}), the instant a node returns. A paused Lux chain is a
recoverable operational event; it is never a divergent chain. This is the crash-fault
contract: \emph{make progress at a majority, or make no unsafe move at all}.

\begin{theorem}[Nova crash-fault safety]
\label{thm:nova}
Under the crash-fault model (validators may halt but do not equivocate), if
\(\mathrm{NovaQuorum}(n)=\lfloor n/2\rfloor+1\) distinct validators must accept a block at a
height before it is Nova-final, then no two conflicting blocks are ever both Nova-final at
that height.
\end{theorem}
\begin{proof}
Let \(B\) and \(B'\) be Nova-final at the same height, certified by acceptor sets \(S,S'\)
with \(|S|,|S'|\ge \lfloor n/2\rfloor+1\). Then
\(|S\cap S'| \ge |S|+|S'|-n \ge 2(\lfloor n/2\rfloor+1)-n \ge 1\), so some validator \(v\in
S\cap S'\) accepted both. Under the crash-fault model a validator never accepts two
conflicting blocks at one height (it accepts, then halts or continues, but does not
equivocate), so \(B=B'\). A network partition cannot manufacture two majorities: at most one
side of any partition can hold \(\lfloor n/2\rfloor+1\) of \(n\).
\end{proof}

Theorem~\ref{thm:nova} is exactly why Nova is safe to \emph{act on locally} but is still
labeled \emph{reorgable}: the crash-fault guarantee does not cover a genuine equivocator, who
could sign for two blocks and split a bare majority. Guarding against that adversary is not
Nova's job --- it is Quasar's.

\section{Quasar --- \(\tfrac{2}{3}\)-stake export finality}
\label{sec:quasar}

Quasar is the \emph{upper} load-bearing tier. It produces a portable, post-quantum
certificate that a strict supermajority of \emph{stake} attested to a block Nova already
accepted. It is the only tier the invariant permits to leave the chain.

\subsection{The stake predicate: one definition, strict \(>\tfrac{2}{3}\)}

The supermajority floor has exactly one definition
(\texttt{config/quorum\_threshold.go}), overflow-safe, shared by both the certificate
verifier and the live-set parameter sizer so the two can never drift:

\begin{lstlisting}[language=Go]
// TwoThirdsStakeFloor(total) = floor(2*total/3), overflow-safe.
// A Quasar certificate is valid iff signed stake is STRICTLY greater.
func TwoThirdsStakeFloor(total uint64) uint64 {
    q, r := total/3, total%3
    floor := 2 * q
    if r == 2 { floor++ }          // floor(2r/3): r in {0,1}->0, r==2->1
    return floor
}
// Equal-stake vote-count form: floor(2n/3)+1.
func EqualStakeSupermajorityThreshold(n int) int {
    if n <= 0 { return 1 }
    return int(TwoThirdsStakeFloor(uint64(n))) + 1
}
\end{lstlisting}

\subsection{The certificate carries its tier}

A \texttt{QuorumCert} is tagged with the rung it certifies, and verification dispatches on
that tag --- the two-tier ladder is not a convention layered over the wire, it \emph{is} the
wire (\texttt{engine/chain/cert.go}):

\begin{lstlisting}[language=Go]
func (c *QuorumCert) VerifyWeighted(v VoteVerifier, stake StakeSource, epochHeight uint64) error {
    if err := c.Verify(v, epochHeight); err != nil { return err }
    switch c.Tier {
    case Nova:   return c.verifyNovaMajority(stake, epochHeight)       // >= floor(n/2)+1 voters
    case Quasar: return c.verifyQuasarSupermajority(stake, epochHeight) // > floor(2/3 * stake)
    default:     return fmt.Errorf("%w: tier=%s", ErrQCUnknownTier, c.Tier)
    }
}

func (c *QuorumCert) verifyQuasarSupermajority(stake StakeSource, epochHeight uint64) error {
    total := stake.TotalStake(epochHeight)          // read at the P-CHAIN EPOCH height
    if total == 0 { return ...fail-closed... }
    var voted uint64
    for i := range c.Votes { voted += stake.Weight(c.Votes[i].NodeID, epochHeight) }
    if voted <= config.TwoThirdsStakeFloor(total) { return ...below supermajority... }
    return nil                                       // STRICT > floor(2/3 * total)
}
\end{lstlisting}

Three details make this correct rather than merely plausible:
\begin{enumerate}[leftmargin=1.4em]
\item \textbf{Strict inequality.} Acceptance is \(\texttt{voted} > \lfloor\tfrac{2}{3}\,
\texttt{total}\rfloor\) (Tendermint \(+\tfrac{2}{3}\)), never \(\ge\).
\item \textbf{Read at the epoch height.} The stake tally is read at the block's P-chain
\emph{epoch} height --- the same height the certificate's set-root and per-voter public keys
are read at --- not the value-chain height, which races ahead of the P-chain and which the
platform VM would reject as unfinalized. This binds every counted vote to the weight the
signer actually held when it signed.
\item \textbf{One floor everywhere.} \texttt{verifyQuasarSupermajority} and the live-set
\(\alpha\)-sizer both call \texttt{config.TwoThirdsStakeFloor}, so the count threshold a node
sizes to can never disagree with the stake predicate a verifier enforces.
\end{enumerate}

\subsection{Trails acceptance; gates only export}

Quasar runs \emph{strictly after} \texttt{VM.Accept}. Each validator, \emph{having already
accepted} block \(B\), signs the canonical finality message for \(B\) (which binds the inner
canonical execution identity and the epoch set-root and excludes the outer proposer-VM
envelope); a collector aggregates signatures bucketed by \((\text{height},\ \text{inner-id})\)
and emits one certificate when the bucket clears the stake floor. Consequently:

\begin{itemize}[leftmargin=1.4em]
\item An absent, late, or never-formed Quasar certificate \textbf{cannot} stall the chain,
revert a block, or change which block Nova accepted. Quasar \emph{describes} an outcome; it
does not \emph{decide} one.
\item What Quasar gates is \emph{export}, not \emph{progress}. A node below its stake
supermajority keeps climbing to Nova and executing locally; it simply does not emit the
export certificate. This is the honest degraded mode, not a failure: local progress
continues, external claims wait.
\item The signed bytes are the unchanged \texttt{CanonicalVoteMessage}, the certificate is
the unchanged \texttt{QuorumCert}, and external verification is the unchanged
\texttt{VerifyWeighted}; bridges, \texttt{0x9999} DEX receipts, and light clients are
unaffected by anything internal to the engine (LP-182, LP-217).
\end{itemize}

\subsection{Post-quantum composition}

The certificate's signature legs are the Lux PQ suite, selected by the operator's cert
profile (LP-217): \texttt{PQ-off} (BLS only), \texttt{PQ-fast} (BLS \(\parallel\) Pulsar),
\texttt{PQ-strict} (BLS \(\parallel\) Pulsar \(\parallel\) Corona, mainnet default),
\texttt{PQ-heavy} (adds Magnetar SLH-DSA). Pulsar contributes Module-LWE threshold signatures
byte-equal to FIPS-204 ML-DSA; Magnetar contributes hash-based SLH-DSA; the Horizon rung
seals the certificate so that it is irreversible even against a quantum adversary. Because
Quasar is post-accept, \textbf{the post-quantum cost is off the critical finality path} --- the
chain has already advanced to Nova when the certificate is being assembled.

\section{The two-tier separation as a safety theorem}
\label{sec:safety}

The value of separating the tiers is that each carries exactly one obligation and the two
compose without a seam.

\subsection{Quasar's Byzantine overlap bound}

The metastable sampler is configured so that two acceptance quorums must overlap in
\emph{more than} \(f\) honest nodes, where \(f=\lfloor (K-1)/3\rfloor\). This is the
BFT-overlap bound enforced in \texttt{config.Parameters.Valid()}:
\[
2\alpha - K \;\ge\; \left\lfloor \tfrac{K-1}{3} \right\rfloor + 1 \;=\; f+1 ,
\]
i.e.\ \(2\alpha - K > f\). Configurations that fail it (e.g.\ \(K{=}3,\alpha{=}2\) posing as
BFT) are rejected with \texttt{ErrAlphaBelowBFTQuorum}; mainnet additionally requires
\(K\ge 11\) so \(f\ge 3\), and a value/PoS chain requires \(K\ge 4\). The node
\emph{fails closed}: it refuses to start a multi-node value chain with non-BFT parameters
rather than run an unsafe quorum.

\begin{theorem}[No two conflicting Quasar certificates]
\label{thm:quasar}
Assume honest validators sign a finality message only for a block they accepted, Nova accepts
at most one block per height (Theorem~\ref{thm:nova}, w.h.p.\ under the sampler margin), and
Byzantine stake is \(< \tfrac{1}{3}\) of total. Then no two Quasar certificates for
conflicting blocks \(B\ne B'\) at the same height can both verify.
\end{theorem}
\begin{proof}
A valid Quasar certificate for \(B\) carries signed stake \(w>\lfloor\tfrac{2}{3}T\rfloor\),
so strictly more than \(\tfrac{2}{3}\) of total stake \(T\). Two such certificates for
\(B\ne B'\) would together require \(> \tfrac{4}{3}T\) signed stake, exceeding \(T\) by
\(>\tfrac{1}{3}T\); by inclusion--exclusion the overlap carries \(>\tfrac{1}{3}T\) of stake
that signed \emph{both}. Since Byzantine stake is \(<\tfrac{1}{3}T\), the overlap contains
honest stake, and an honest validator signs only what it accepted, so it signed at most one
of \(B,B'\) --- contradiction. Hence at most one conflicting certificate verifies.
\end{proof}

\subsection{Why the reduction is clean}

Quasar adds \emph{no} new safety assumption and \emph{no} liveness dependency beyond what the
legacy \(\tfrac{2}{3}\)-stake certificate already carried. Its safety reduces to Nova's
single-accept-per-height (Theorem~\ref{thm:nova}) plus the same \(\ge\tfrac{2}{3}\)-honest-
stake assumption; nothing waits on it. The composition is therefore:
\[
\underbrace{\text{Nova (majority)}}_{\text{liveness to }n/2,\ \text{crash-safe local accept}}
\;\prec\;
\underbrace{\text{Quasar (}\tfrac{2}{3}\text{ stake)}}_{\text{Byzantine-safe export}}
\;\prec\;
\underbrace{\text{Horizon (PQ seal)}}_{\text{quantum-irreversible}} .
\]
Each rung is independently complete; each authorizes strictly more than the last; and the
invariant --- nothing below Quasar exports --- is a single comparison, \(\texttt{finality}
\ge \texttt{Quasar}\), checked at every export boundary (\texttt{Finality.AuthorizesExport}).

\subsection{Where Lux sits among the thresholds}

\begin{center}
\begin{tabular}{@{}l l l@{}}
\toprule
\textbf{Family} & \textbf{Safety threshold} & \textbf{Finality} \\
\midrule
\textbf{Lux Nova (local)}   & majority \(\lfloor n/2\rfloor+1\), crash-safe & sub-second, reorgable \\
\textbf{Lux Quasar (export)}& strict \(>\tfrac{2}{3}\) stake, Byzantine     & deterministic PQ certificate \\
Classical BFT (Tendermint, HotStuff) & \(\tfrac{1}{3}\) of a fixed committee & instant, but \(O(n^2)\), does not scale \\
Nakamoto (Bitcoin)          & \(>\!50\%\) hashpower                          & probabilistic, minutes \\
\bottomrule
\end{tabular}
\end{center}

Lux keeps the metastable sampler's scale (a fixed small sample finalizes among millions of
validators) and \emph{adds} a deterministic post-quantum export certificate --- without
paying BFT's quadratic message cost or Nakamoto's latency. The full cross-system comparison is
LP-306~\cite{lp306}.

\section{ZAP-native transport --- no gRPC, no round machine}
\label{sec:zap}

Nova's messages ride the ZAP wire (\texttt{luxfi/zap v0.8.1}, a direct dependency of
\texttt{consensus v1.36.0}; LP-022 / LP-200--204). ZAP is the single Lux plugin transport:
the engine's \texttt{Transport[ID]} interface (\texttt{engine/chain/engine.go}) is satisfied
by the ZAP plugin, and there is \textbf{no gRPC plugin server, no Tendermint p2p, and no
BFT round/gossip state machine} anywhere on the acceptance path. This matters for the ladder
in three ways:

\begin{itemize}[leftmargin=1.4em]
\item \textbf{Sampling, not broadcast.} Photon queries a \(k\)-subset per round; the transport
carries \(O(k)\) point-to-point polls, not an \(O(n^2)\) all-to-all vote. The transport cost
of a round is independent of \(n\).
\item \textbf{Post-quantum by construction.} ZAP frames carry the Lux PQ suite, so a Quasar
attestation is a signed ZAP message; the certificate is an aggregate of ZAP-carried
signatures. No transport translation layer sits between the signature and the wire.
\item \textbf{No transport-level view or leader.} Because there is no BFT round machine, there
is no view-change, no leader-election message class, and no lock/precommit traffic to carry ---
the wire has nothing to say about ``rounds'' because the decider has no rounds.
\end{itemize}

\section{Size-invariance: the same rule from \(n=1\) to \(n=10^6\)}
\label{sec:scale}

\(K\) is the \textbf{sample} size, not the validator-set size --- the load-bearing
distinction. For \(n\le K\) a node polls the whole set (\(k=n\)); for \(n>K\) it polls a fixed
\(k\)-subset each round. The \emph{rule} --- repeated \(k\)-sampling, \(\alpha\)-of-\(k\)
threshold, \(\beta\)-confidence, accept at the Nova majority, export at the Quasar
supermajority --- is identical at every scale, and the per-round cost stays \(O(k)\) per node
regardless of \(n\).

\begin{center}
\begin{tabular}{@{}l r r r l@{}}
\toprule
\textbf{Deployment} & \(K\) & \(\alpha\) & \(\beta\) & \textbf{Source} \\
\midrule
Primary-network fleet (live)   & 5  & 4  & 2  & \texttt{node/chains/quorum.go} live sizing \\
\texttt{DefaultParams}         & 20 & 14 & 14 & \texttt{config/config.go} \\
\texttt{MainnetParams} (floor) & 21 & 15 & --- & \texttt{config/config.go} \\
Testnet                        & 11 & 8  & --- & \texttt{config/config.go} \\
\bottomrule
\end{tabular}
\end{center}

At the live operating point \(K{=}5, \alpha{=}4\) the sampler's per-round agreement bar is
\(\alpha/K = 80\%\), comfortably clearing the BFT-overlap bound \(2\alpha-K = 3 \ge f+1 = 2\).
Note the two thresholds this paper separates are visible here as \emph{distinct knobs}: the
sampler's per-round \(\alpha\) (how many of the \(k\) polled must agree to score a round) is
not the Nova acceptance count (a majority of the set) and is not the Quasar export floor
(\(>\tfrac{2}{3}\) stake). What an operator changes as \(n\) grows is only which small \(k\)
to pin; the ladder never changes.

This is exactly the property that a ``\(\tfrac{2}{3}\)-of-\emph{all}-stake certificate
decides'' model destroys: collecting a supermajority of \emph{every} validator's signature per
height is \(O(n)\) work on the critical path. Nova restores size-invariance by keeping the
decider a fixed-sample metastable process, and keeps the \(\tfrac{2}{3}\)-stake artifact ---
but moves it \emph{off} the critical path, onto the Quasar export rung.

\section{Empirical evaluation}
\label{sec:eval}

\paragraph{Environments.} Two environments, named once here and referenced by tag:
\begin{itemize}[leftmargin=1.4em]
\item \textbf{Cluster} (\Mnow): the live Lux \emph{testnet}, \num{5} validators
\texttt{luxd-0..4} running \texttt{ghcr.io/luxfi/node:v1.36.0} (\texttt{consensus v1.36.0},
\texttt{evm v1.104.8}), namespace \texttt{lux-testnet} on \texttt{do-sfo3-lux-k8s}, C-Chain
id \num{96368}, internal RPC \texttt{/v1/bc/C/rpc} port \num{9640}. Measured 2026-07-10.
\item \textbf{Host} (\Mnow): Apple M1 Max, 10 cores, \SI{64}{GB}, macOS 26.5.1, \texttt{go
1.26.4}, \texttt{arm64}. Used for the in-process consensus-engine microbenchmarks.
\end{itemize}

\subsection{Coherence under a node-drop --- the two-tier property, live}
\label{sec:eval:drop}

The one property the two-tier ladder exists to deliver is: \emph{a validator falls behind, the
chain keeps finalizing locally, the validator rejoins, and no fork ever appears}. We drove a
five-validator testnet under steady load (24 in-cluster senders) and scaled the validator set
\(5 \to 4 \to 3 \to 4 \to 5\), measuring at each phase the Nova (latest) height advance and
\emph{common-height coherence} --- all live nodes must hold the \emph{same} block hash at a
confirmed height (height skew is not a fork; a hash disagreement is). Harness:
\texttt{bench/vs-ava}\slash\ (\texttt{nova\_proof.sh}, \texttt{nova\_loadgen.py}).

\begin{center}
\begin{tabular}{@{}l r c c l@{}}
\toprule
\textbf{Phase} & \textbf{Nova advance} & \textbf{Coherent} & \textbf{Tier} & \textbf{Note} \\
 & (blocks / \SI{70}{s}) & (1 hash @ conf.) & & \\
\midrule
\(5/5\) baseline        & \(+20\) & yes & Nova + Quasar & both tiers advance \\
\(4/5\) drop \texttt{luxd-4} & \(+8\)  & yes & Nova + Quasar & \(4/5 \ge \tfrac{2}{3}\): export still forms \\
\(\mathbf{3/5}\) drop \texttt{luxd-3} & \(\mathbf{+6}\)  & \textbf{yes} & \textbf{Nova only} & \textbf{majority holds; \(<\tfrac{2}{3}\) so export honestly pauses} \\
\(4/5\) restore \texttt{luxd-3} & (rejoin) & yes & Nova + Quasar & \texttt{luxd-3} caught up \(150\!\to\!161\), converged \\
\(5/5\) restore \texttt{luxd-4} & (rejoin) & yes & Nova + Quasar & full set restored, coherent \\
\bottomrule
\end{tabular}
\end{center}

\noindent All phases held \texttt{distinct@conf = 1} (a single common-height hash across live
nodes) and \texttt{clamp\_ok} (the exported/finalized tag never exceeded the latest height).
The load-bearing row is \(\mathbf{3/5}\): the majority tier keeps producing and stays coherent
at exactly the point a \(\tfrac{2}{3}\)-only engine freezes (\(\lceil 2\cdot5/3\rceil=4>3\)),
while the export tier \emph{honestly pauses} because no \(>\tfrac{2}{3}\)-stake certificate can
form on three of five. The dropped validator rejoined and converged to the common chain
(observed \texttt{luxd-3}: \(150\to161\)). This is Nova (majority, \S\ref{sec:nova}) and Quasar
(\(\tfrac{2}{3}\)-stake, \S\ref{sec:quasar}) doing precisely their separate jobs, live. \Mnow

\subsection{Live C-Chain throughput --- the ceiling is the fee/gas-target, not consensus}
\label{sec:eval:tput}

To locate the live throughput ceiling we ran two in-cluster loads against the same
five-validator testnet over \(\sim\!\SI{90}{s}\) windows: a \emph{closed-loop} run (each
account advances one nonce only after the prior tx is accepted;
\texttt{bench/vs-ava/tp\_loadgen.py}) and an \emph{open-loop, saturating} run (a deep
per-account nonce pipeline floods the mempool so the block builder is the only throttle;
\texttt{tp\_loadgen\_open.py}). The result is unambiguous under \emph{both} loads: the chain
mines a low transaction rate \emph{at low block fullness with the base fee pinned at its floor}
--- so the binding constraint is the EVM's dynamic-fee / gas-target throttle, not consensus,
which stayed fully coherent.

\begin{center}
\begin{tabular}{@{}l r r l@{}}
\toprule
\textbf{Metric} & \textbf{closed-loop} & \textbf{open-loop (sat.)} & \textbf{Tag} \\
\midrule
window / blocks               & \SI{91}{s} / 18 & \SI{90}{s} / 21 & \Mnow \\
mined throughput (tx/s)       & \num{1.81}      & \num{5.72}      & \Mnow \\
mean block interval           & \SI{4.38}{s}    & \SI{4.15}{s}    & \Mnow \\
mean tx / block               & \num{9.17}      & \num{24.6} (peak \textbf{106}) & \Mnow \\
\textbf{block fullness}       & \(\mathbf{1.60\%}\) & \(\mathbf{4.31\%}\) mean, \(\mathbf{18.55\%}\) peak & \Mnow \\
base fee                      & \SI{25}{gwei} floor & \SI{25}{gwei} floor & \Mnow \\
cross-node coherence          & \textbf{yes} (all 5 = \texttt{0x5960dc41\dots}) & \textbf{yes} (all 5 = \texttt{0x9fb2b542\dots}) & \Mnow \\
\bottomrule
\end{tabular}
\end{center}

Even when \emph{saturated} (open-loop), blocks peak at only \(\sim\!19\%\) full with the base
fee still at its \SI{25}{gwei} floor, and consensus stays coherent --- so the chain is
nowhere near execution-bound or consensus-bound; blocks are throttled by the coreth/evm
gas-target and block-cadence policy, not by the sampler and not by the EVM's execution
capacity. Consensus was coherent and idle throughout (identical hashes across all five nodes).
This reproduces, on v1.36, the prior in-cluster finding that the C-Chain's live mined ceiling
is a \emph{fee-policy} ceiling: consensus absorbs offered load coherently, and the EVM
fee/gas-target caps what is mined~\cite{luxquasarbench}~(\Mprior). It is the honest frame for
the head-to-head against \texttt{avalanchego} in \S\ref{sec:eval:exec}: both C-Chains run the
same sequential, fee-throttled, geth-derived EVM, so their \emph{live} mined rates are
governed by the same throttle --- the Lux advantage at v1.36 is the consensus-liveness axis
(\S\ref{sec:eval:drop}), not raw live TPS.

\subsection{Relationship to execution --- Block-STM is downstream and (today) sequential}
\label{sec:eval:exec}

Execution is orthogonal to the decider and \emph{downstream} of it: Nova accepts, then the EVM
executes the accepted block. On the live C-Chain that execution is \textbf{100\%\ sequential
coreth-lineage EVM} --- the parallel GPU-native Block-STM executor is built and benchmarked but
\emph{not wired into the production path} (\texttt{RegisterExecutor} has zero callers;
\texttt{DefaultExecutor()} is a no-op fallback;
\texttt{evm/core/state\_processor.go} falls through to the sequential loop). This is the same
execution model \texttt{avalanchego}'s C-Chain runs, and the honest reason the two chains'
\emph{live} throughput is comparable. The measured parallel-executor headroom (CPU
copy-bound today; projected on GPU) and the standalone \texttt{avalanchego} head-to-head are
reported in the companion benchmark study (\texttt{bench/vs-ava}) and the GPU-native execution
paper~\cite{gpubtps}; consensus does not gate any of it.

\subsection{Consensus engine bookkeeping is not the bottleneck}
\label{sec:eval:engine}

The in-process chain-engine microbenchmarks (\texttt{engine/chain/benchmark\_test.go}, Host
\Mnow, median of 5) confirm that block/vote/finalization \emph{bookkeeping} runs orders of
magnitude above anything the fee-throttled EVM asks of it:

\begin{center}
\begin{tabular}{@{}l r r@{}}
\toprule
\textbf{Operation} & \textbf{median} & \textbf{rate} \\
\midrule
finalization check     & \SI{19}{ns}   & \(\sim\)\num{52} M/s \\
vote processing        & \SI{424}{ns}  & \(\sim\)\num{2.4} M/s \\
block acceptance       & \SI{251}{ns}  & \(\sim\)\num{4.0} M/s \\
block ingestion        & \SI{1.2}{\micro s} & \(\sim\)\num{0.83} M/s \\
\bottomrule
\end{tabular}
\end{center}

\noindent Every engine operation is sub-microsecond to low-microsecond --- six-plus orders of
magnitude above the \(\sim\!\num{1.8}\) tx/s and \(\sim\!\num{0.2}\) blk/s the fee-throttled
live chain produces (\S\ref{sec:eval:tput}). The engine's data-structure throughput is
therefore never the live ceiling. (This microbenchmark exercises the engine's bookkeeping
harness, not the full Nova quorum decision, and is reported as a lower bound on engine
headroom, not as the Nova finality rate. Raw:
\texttt{bench/vs-ava/results/engine\_bench\_raw.txt}.)

\section{Terminology reconciliation and relationship to prior LPs}
\label{sec:reconcile}

``Nova'' has been used three ways in the Lux corpus; v1.36 fixes the meaning to the one in the
shipped code, and this section states the reconciliation so no reader conflates them.

\begin{itemize}[leftmargin=1.4em]
\item \textbf{v1.36 Nova (this paper, canonical).} A \emph{rung} on the finality ladder ---
majority local-acceptance --- and, by extension, the \emph{linear-chain consensus mode} that
climbs that ladder (\texttt{protocol/nova}, \texttt{engine/chain/finality.go}). This is the
meaning in \texttt{consensus v1.36.0} and in LP-134's taxonomy (Nova = linear mode; Nebula =
DAG mode).
\item \textbf{LP-020 ``Nova''} named the linear-chain \emph{engine} (as opposed to Nebula the
DAG engine); its component vocabulary (Photon / Wave / Focus / Prism / Horizon / Flare) and its
QuasarCert wire survive here \emph{as the attestation layer}. Consistent with this paper once
one reads ``Nova'' as ``linear-chain mode.''
\item \textbf{The earlier ``supernova-burst'' DAG Nova} (a chapter of
\texttt{lux-consensus-engines}) described a DAG-frontier burst finalizer. That is a
\emph{Nebula}-family (DAG-mode) design; it is \emph{not} the v1.36 linear-chain two-tier
ladder specified here, and where the two disagree on what the ``Nova'' rung means, this paper
and LP-305 govern for the shipped linear-chain C-Chain.
\end{itemize}

\noindent Relationship to the specifications:
\begin{itemize}[leftmargin=1.4em]
\item \textbf{LP-305}~\cite{lp305} is the normative spec of record for this ladder; this paper
is its rigorous, measured companion.
\item \textbf{LP-306}~\cite{lp306} is the cross-system landscape (Lux vs Tendermint /
Avalanche / Nakamoto / Gasper / Sui / Aptos / Solana).
\item \textbf{LP-105} remains the authoritative \emph{lexicon} (public vocabulary and internal
names); its finality \emph{state-machine} description predates Nova and is refined by this
ladder --- read the Nova/Quasar tiers here for the mechanism, LP-105 for the names.
\item \textbf{LP-010 / LP-210} specify Block-STM (the execution layer, \S\ref{sec:eval:exec});
\textbf{LP-303} the bounded re-poll / catch-up that heals a paused chain; \textbf{LP-217} the
cert profile modes (\S\ref{sec:quasar}).
\end{itemize}

\section{Conclusion}

Nova's contribution is not a new voting rule; it is a \emph{decomposition}. By splitting
finality into a majority \emph{local-accept} tier and a \(\tfrac{2}{3}\)-stake \emph{export}
tier --- and forbidding, by a one-line invariant, any export below the supermajority --- Lux
v1.36 gets crash-fault liveness (three of five keep producing) \emph{and} Byzantine-safe
external finality (nothing leaves the chain without \(>\tfrac{2}{3}\) of stake), without the
false choice a single threshold forces. The decider stays a fixed-sample metastable process,
so the rule is size-invariant from one validator to a million; the \(\tfrac{2}{3}\)-stake
certificate stays post-quantum and portable, but rides the upper rung, off the critical path.
The live testnet shows the two tiers doing their separate jobs: coherence held through a
\(5/5\to4/5\to3/5\) drop while the export tier honestly paused, and the live throughput ceiling
sat in the EVM fee policy, not in consensus. The ladder is the spec; the code is the proof.

\begin{thebibliography}{9}
\bibitem{lp305} Lux Industries. \emph{LP-305: Lux Consensus --- Canonical Specification (Nova
\(\cdot\) Quasar \(\cdot\) Block-STM).} \url{https://github.com/luxfi/lps}.
\bibitem{lp306} Lux Industries. \emph{LP-306: Consensus Comparison --- Lux (Nova/Quasar) vs
Tendermint, Avalanche, Nakamoto, Gasper, Sui, Aptos, Solana.}
\url{https://github.com/luxfi/lps}.
\bibitem{lp303} Lux Industries. \emph{LP-303: Bounded Re-poll and Runtime Ancestor Catch-up.}
\url{https://github.com/luxfi/lps}.
\bibitem{gpubtps} Lux Industries. \emph{The Path to One Billion Trades per Second: GPU-Native
Execution on Lux.} Lux Papers, 2026.
\bibitem{luxquasarbench} Lux Industries. \emph{Lux Quasar Benchmarks and In-Cluster C-Chain
Throughput Measurements.} Lux Papers, 2026.
\bibitem{avalanche2020} Team Rocket et al. \emph{Scalable and Probabilistic Leaderless BFT
Consensus through Metastability.} 2020. \url{https://arxiv.org/abs/1906.08936}.
\bibitem{blockstm2022} Gelashvili et al. \emph{Block-STM: Scaling Blockchain Execution by
Turning Ordering Curse to a Performance Blessing.} 2022. \url{https://arxiv.org/abs/2203.06871}.
\end{thebibliography}

\end{document}
